\documentclass[11pt,a4paper]{article}
\usepackage[T1]{fontenc}
\usepackage{isabelle,isabellesym}

% further packages required for unusual symbols (see also
% isabellesym.sty), use only when needed

\usepackage{amssymb}
  %for \<leadsto>, \<box>, \<diamond>, \<sqsupset>, \<mho>, \<Join>,
  %\<lhd>, \<lesssim>, \<greatersim>, \<lessapprox>, \<greaterapprox>,
  %\<triangleq>, \<yen>, \<lozenge>

%\usepackage{eurosym}
  %for \<euro>

\usepackage[only,bigsqcap,bigparallel,fatsemi,interleave,sslash]{stmaryrd}
  %for \<Sqinter>, \<Parallel>, \<Zsemi>, \<Parallel>, \<sslash>

%\usepackage{eufrak}
  %for \<AA> ... \<ZZ>, \<aa> ... \<zz> (also included in amssymb)

\usepackage{textcomp}
  %for \<onequarter>, \<onehalf>, \<threequarters>, \<degree>, \<cent>,
  %\<currency>

\usepackage{wasysym}
\usepackage{pifont}

% this should be the last package used
\usepackage{pdfsetup}

% urls in roman style, theory text in math-similar italics
\urlstyle{rm}
\isabellestyle{it}

% for uniform font size
%\renewcommand{\isastyle}{\isastyleminor}


\begin{document}

\title{Completeness of $Q_0$}
\author{Asta Halkjær Boserup\and Jonathan Juli\'an Huerta y Munive\and Anders Schlichtkrul}
\maketitle

\begin{abstract}
  We formalize the completeness of Peter B. Andrews' system $Q_0$ for Higher-Order Logic 
  following his textbook ``An Introduction to Mathematical Logic and Type Theory: To Truth Through Proof''. 
  The formalization builds on Diaz's formalization of $Q_0$'s syntax, semantics, soundness and consistency. 
  The completeness is with respect to general models. We prove completeness by introducing an abstract 
  consistency property for $Q_0$ and using the framework for abstract consistency properties by From and 
  Schlichtkrull to get a model existence theorem for $Q_0$. In our application of the framework we also adapt
  proofs from Andrews' book to this context.
\end{abstract}

\tableofcontents

% sane default for proof documents
\parindent 0pt\parskip 0.5ex

% generated text of all theories
\input{session}

\nocite{*}

% optional bibliography
\bibliographystyle{abbrv}
\bibliography{root}

\end{document}

%%% Local Variables:
%%% mode: latex
%%% TeX-master: t
%%% End:
